\section*{Capítulo \ref{ch:g11:absorbance} --- Color y absorbancia}

\begin{solution}{exo:g11:absorbance:1}
Si absorbe el naranja, se ve azul, el color enfrentado al naranja en el
círculo. Si absorbe el violeta, se ve amarilla.
\end{solution}

\begin{solution}{exo:g11:absorbance:2}
Una disolución verde absorbe sobre todo la luz roja, el color enfrentado
al verde en el círculo: de unos \qty{620}{nm} a \qty{750}{nm}.
\end{solution}

\begin{solution}{exo:g11:absorbance:3}
Colorante azul: $\lambda_{\max} = \qty{629}{nm}$, $A = 0.82$. Colorante
amarillo: $\lambda_{\max} = \qty{427}{nm}$, $A = 0.795$, alrededor de
0.80.
\end{solution}

\begin{solution}{exo:g11:absorbance:4}
$A$ es proporcional a la concentración: el triple da $A = 0.90$; la
mitad da $A = 0.15$.
\end{solution}

\begin{solution}{exo:g11:absorbance:5}
El blanco es una cubeta llena del \omterm{def:g6:solutions-and-solubility:solute}{disolvente} puro. Hacer $A = 0$ con él
elimina lo que absorben la cubeta, el \omterm{def:g6:solutions-and-solubility:solute}{disolvente} y el instrumento, de
modo que las lecturas posteriores se deben solo al \omterm{def:g6:solutions-and-solubility:solute}{soluto}.
\end{solution}

\begin{solution}{exo:g11:absorbance:6}
$C = 0.41 / 0.164 = \qty{2.5}{mg/L}$.
\end{solution}

\begin{solution}{exo:g11:absorbance:7}
A \qty{629}{nm} la \omterm{def:g11:absorbance:absorbance}{absorbancia} es máxima, así que la medida es la más
sensible (concentraciones pequeñas dan todavía \omterm{def:g11:absorbance:absorbance}{absorbancias} legibles); y
en lo alto de la banda la curva es plana, así que un pequeño error en la
longitud de onda apenas cambia la lectura. A \qty{550}{nm} la \omterm{def:g11:absorbance:absorbance}{absorbancia}
es pequeña y cambia deprisa con la longitud de onda.
\end{solution}

\begin{solution}{exo:g11:absorbance:8}
Diluida: $C = 0.48 / 0.164 = \qty{2.9}{mg/L}$; la bebida:
$5 \times 2.9 = \qty{15}{mg/L}$ (con más precisión, $14.6$).
\end{solution}

\begin{solution}{exo:g11:absorbance:9}
$a = A / (\ell\, C_m) = 0.795 / (1.00 \times 0.0150) = \qty{53.0}{L/(g.cm)}$;
$\varepsilon = a \times M = 53.0 \times 534.4 = \qty{2.83e4}{L/(mol.cm)}$.
\end{solution}

\begin{solution}{exo:g11:absorbance:10}
Por encima de unos \qty{520}{nm}. A \qty{629}{nm} el colorante amarillo
no absorbe nada, así que toda la \omterm{def:g11:absorbance:absorbance}{absorbancia} a esa longitud de onda
procede del colorante azul.
\end{solution}

\begin{solution}{exo:g11:absorbance:11}
Se duplica: $A = \varepsilon \ell c$ es proporcional a $\ell$; la luz
atraviesa el doble de disolución y, por tanto, el doble de \omterm{def:g7:atoms-and-molecules:molecule}{moléculas}
absorbentes.
\end{solution}

\begin{solution}{exo:g11:absorbance:12}
$A/C$: $0.164$, $0.162$, $0.115$, $0.0675$. El cociente solo es
constante para los dos primeros: por encima de unos \qty{10}{mg/L}
(\omterm{def:g11:absorbance:absorbance}{absorbancias} superiores a 1.6, más o menos) la ley deja de cumplirse.
Solo se pueden usar los patrones de 5 y \qty{10}{mg/L} (y otros más
bajos); una desconocida que marca 2.5 hay que diluirla, por ejemplo
cinco veces, y volver a medirla.
\end{solution}

\begin{solution}{exo:g11:absorbance:13}
Colorante azul, solo a partir de \qty{629}{nm}:
$C = 0.574 / 0.164 = \qty{3.5}{mg/L}$. Colorante amarillo, a partir de
\qty{427}{nm}: $C = 0.53 / 0.0530 = \qty{10}{mg/L}$.
\end{solution}

\begin{solution}{exo:g11:absorbance:14}
Obtiene $0.368 / 0.164 = \qty{2.24}{mg/L}$. La \omterm{def:g11:absorbance:absorbance}{absorbancia} real es
$0.368 - 0.040 = 0.328$, así que la concentración real es
$0.328 / 0.164 = \qty{2.00}{mg/L}$. Se pasa en un \qty{12}{\percent}.
\end{solution}

\begin{solution}{exo:g11:absorbance:15}
$10 \times 60 = \qty{600}{mg}$ al día, en $600 / 20 = \qty{30}{L}$ de
bebida. Nadie bebe treinta litros al día: la bebida sola no supone un
riesgo realista, aunque el colorante también puede venir de otros
alimentos.
\end{solution}

\begin{solution}{pb:g11:absorbance:1}
\textbf{1.} La rojo anaranjada, el color enfrentado al azul en el
círculo.

\textbf{2.} $\lambda_{\max} = \qty{629}{nm}$.

\textbf{3.} A \qty{450}{nm} la luz es azul: el colorante la deja pasar,
y por eso precisamente la bebida se ve azul.

\textbf{4.} A \qty{629}{nm}.

\textbf{5.} $\qty{50.0}{mg} / \qty{0.5000}{L} = \qty{100}{mg/L}$.

\textbf{6.} \omterm{def:g10:concentration-and-dilution:dilution}{Factor de dilución} $100 / 3.0 = 33.3$, así que
$V_0 = 100.0 / 33.3 = \qty{3.00}{mL}$: una pipeta graduada (o una
bureta) y un matraz aforado de \qty{100.0}{mL}.

\textbf{7.} $A/C$ = 0.168, 0.163, 0.165, 0.163, 0.165: todos cerca de
0.164, así que $A \approx 0.164\, C$, una recta que pasa por el origen.

\textbf{8.} Una disolución sin colorante ($C = 0$) es el blanco, con el
que se hace $A = 0$.

\textbf{9.} $a = \qty{0.164}{L/mg}$ por centímetro, es decir,
\qty{164}{L/(g.cm)}; $\varepsilon = 164 \times 792.9 =
\qty{1.30e5}{L/(mol.cm)}$.

\textbf{10.} $F = 50.0 / 25.0 = 2$.

\textbf{11.} $C = 0.590 / 0.164 = \qty{3.60}{mg/L}$.

\textbf{12.} $2 \times 3.60 = \qty{7.20}{mg/L}$.

\textbf{13.} Alrededor de $2 \times 0.590 = 1.18$, por encima del patrón
más alto (0.823): la lectura quedaría fuera del intervalo calibrado,
donde la ley puede fallar. Al diluir, queda entre los patrones.

\textbf{14.} $\qty{7.20}{mg/L} \times \qty{0.500}{L} = \qty{3.60}{mg}$.

\textbf{15.} $n = \num{3.60e-3} / 792.9 = \qty{4.54e-6}{mol}$.

\textbf{16.} $6 \times 30 = \qty{180}{mg}$ al día.

\textbf{17.} $180 / 3.60 = 50$ botellas.

\textbf{18.} $50 \times 0.500 = \qty{25}{L}$ en un día: imposible de
beber. El colorante de esta bebida no puede, por sí solo, acercar a un
niño al límite, aunque los colorantes de varios alimentos se suman.

\textbf{19.} $6 \times 70 = \qty{420}{mg}$, es decir,
$420 / 3.60 \approx
117$ botellas.

\textbf{20.} Unas 50 botellas de \qty{500}{mL}, \qty{25}{L} de bebida,
en un solo día.
\end{solution}
